% ==STYLE-BEGIN==
\documentclass[UTF8,11pt]{ctexart}
\usepackage{amsmath,amssymb,amsthm}
\usepackage[a4paper,margin=2.4cm,top=2.3cm,bottom=2.3cm]{geometry}
\usepackage{fancyhdr}
\usepackage{enumitem}
\usepackage{hyperref}
\hypersetup{hidelinks}

% ---- 版式 ----
\linespread{1.15}
\setlength{\parskip}{0.35em}
\setlength{\parindent}{1.6em}
\setlist{nosep,leftmargin=2em}

% ---- 页眉 ----
\pagestyle{fancy}
\fancyhf{}
\fancyhead[L]{\small\coursename}
\fancyhead[R]{\small\thepage}
\renewcommand{\headrulewidth}{0.4pt}

% ---- 定理环境：中文用正体，标题加粗，按小节共用一条流水号 ----
\newtheoremstyle{cnthm}{0.7em}{0.7em}{\normalfont}{}{\bfseries}{}{0.5em}{}
\theoremstyle{cnthm}
\newtheorem{thm}{定理}[section]
\newtheorem{lem}[thm]{引理}
\newtheorem{prop}[thm]{命题}
\newtheorem{cor}[thm]{推论}
\newtheorem{defn}[thm]{定义}
\newtheorem{rem}[thm]{注}
\newtheorem{exa}[thm]{例}
\newtheorem{conj}[thm]{猜想}
\renewcommand{\proofname}{\bfseries 证明}
\providecommand{\coursename}{}
% ==STYLE-END==
\renewcommand{\coursename}{数学分析\,·\,1019}
\title{数学分析 1019}
\date{}
\begin{document}
\maketitle

\section{连续函数}

\begin{defn}[连续]
$f(x)$ 在 $x_0$ 处连续是指
\[
\lim_{x\to x_0}f(x)=f(x_0).
\]
\end{defn}

\begin{rem}
$\lim\limits_{x\to x_0}f(x)$ 存在 $\iff\lim\limits_{x\to x_0^-}f(x)=\lim\limits_{x\to x_0^+}f(x)$。由上，$f(x)$ 在 $x_0$ 处连续 $\iff$
\[
\lim_{x\to x_0^-}f(x)=\lim_{x\to x_0^+}f(x)=f(x_0).
\]
\end{rem}

\begin{defn}[左连续]
若 $\lim\limits_{x\to x_0^-}f(x)=f(x_0)$，则称 $f$ 在 $x_0$ 处左连续。
\end{defn}

由上，可以按 $f$ 在 $x_0$ 处左、右连续的情形对连续性分类（图略）。

\section{连续函数的性质}

\begin{thm}[四则运算]
若 $f(x)$ 和 $g(x)$ 在 $x=x_0$ 处连续，则 $f(x)\pm g(x)$、$f(x)\cdot g(x)$ 连续；若 $g(x_0)\neq 0$，则 $\dfrac{f}{g}$ 在 $x=x_0$ 处连续。
\end{thm}

\begin{thm}[复合]
若 $g(x)$ 在 $x_0$ 处连续且 $g(x_0)=u_0$，$f(u)$ 在 $u=u_0$ 处连续，则 $f(g(x))$ 在 $x=x_0$ 处连续。
\end{thm}

于是自然要问：是否总有
\[
\lim_{x\to x_0}f(g(x))=f(g(x_0))\,?
\]

\begin{rem}[复合的极限]
若 $\lim\limits_{u\to u_0}f(u)=A$，$\lim\limits_{x\to x_0}g(x)=u_0$，且存在 $x_0$ 的去心邻域使 $g(x)\neq u_0$，则 $\lim\limits_{x\to x_0}f(g(x))=A$。当 $A=f(u_0)$ 时，去心邻域这一条件可以去掉。
\end{rem}

回到上面的问题，有
\[
A=f(u_0)=f(g(x_0))=f\Bigl(\lim_{x\to x_0}g(x)\Bigr)=\lim_{x\to x_0}f(g(x)),
\]
最后一步正是上面的注（取 $A=f(u_0)$，条件可去）。

\begin{rem}
如果 $g(x)$ 在 $x_0$ 附近趋近于 $u_0$，且 $f$ 在 $u_0$ 处连续，那么 $f$ 和 $\lim$ 可以交换。

如果 $f,g$ 在各自的定义域上都连续，那么对复合函数定义域内任何 $x_0$，$f$ 与 $\lim\limits_{x\to x_0}$ 都可以交换。
\end{rem}

\begin{thm}[初等函数]
初等函数在其定义域内连续。
\end{thm}

\section{介值定理}

\begin{lem}[闭区间套定理]
若 $[a_{n+1},b_{n+1}]\subseteq[a_n,b_n]$ 且 $\lim\limits_{n\to+\infty}(b_n-a_n)=0$，则存在唯一的 $c$ 属于所有 $[a_n,b_n]$。
\end{lem}

\begin{thm}[介值定理]
$f:[a,b]\to\mathbb{R}$ 是连续函数，如果 $f(a)<0$ 且 $f(b)>0$，则 $\exists c\in(a,b)$ 满足 $f(c)=0$。
\end{thm}

\begin{proof}
取 $c_1=\frac{a+b}{2}$。若 $f(c_1)=0$，则已成立；若 $f(c_1)>0$，则取 $[a_1,b_1]=[a,c_1]$；若 $f(c_1)<0$，则取 $[a_1,b_1]=[c_1,b]$。可知 $b_1-a_1=\frac{b-a}{2}$。以此类推，如果过程中点始终不为 $0$，则构造出
\[
[a_{n+1},b_{n+1}]\subseteq[a_n,b_n],
\]
且 $b_n-a_n=\frac{1}{2^n}(b-a)$，即
\[
\lim_{n\to+\infty}(b_n-a_n)=0.
\]
由此，$\exists!\,c\in\bigcap\limits_{n=1}^{+\infty}[a_n,b_n]$，且由构造有 $f(a_n)<0$、$f(b_n)>0$，
\[
\lim_{n\to+\infty}a_n=\lim_{n\to+\infty}b_n=c.
\]
易知
\[
f(c)=f\Bigl(\lim_{n\to+\infty}a_n\Bigr)=\lim_{n\to+\infty}f(a_n)\le 0,
\]
同理有 $f(c)\ge 0$，即得 $f(c)=0$。
\end{proof}

\begin{thm}[一般形式的介值定理]
$f:[a,b]\to\mathbb{R}$ 连续，如果 $f(a)<\xi$、$f(b)>\xi$，则 $\exists c\in(a,b)$ s.t. $f(c)=\xi$。
\end{thm}

\begin{rem}
\begin{enumerate}
\item 结论中的 $c$ 取在开区间 $(a,b)$ 内：端点显然不满足。
\item $f:(a,b)\to\mathbb{R}$，如果 $c,d\in(a,b)$、$c<d$ 且 $f(c)<\xi$、$f(d)>\xi$，则 $\exists e\in(c,d)$ s.t. $f(e)=\xi$。（$f$ 在 $(a,b)$ 上连续，显然 $f:[c,d]\to\mathbb{R}$ 连续。）
\item $f:I\to\mathbb{R}$，$I$ 是区间，$f$ 在 $I$ 上连续，$\forall a,b\in I$，则 $f$ 可以取到 $f(a)$ 与 $f(b)$ 中间的所有值，即
\[
\forall\eta\in\bigl[\min\{f(a),f(b)\},\max\{f(a),f(b)\}\bigr],\ \exists x_0\in I\ \text{s.t.}\ f(x_0)=\eta.
\]
进一步有：$\forall\xi_1,\xi_2\in\operatorname{Im}f$（不妨设 $\xi_1\le\xi_2$），$[\xi_1,\xi_2]\subseteq\operatorname{Im}f$。
\end{enumerate}
\end{rem}

问题：如果 $I$ 为有界闭区间，$f$ 连续，那么 $\operatorname{Im}f$ 是什么？

$\operatorname{Im}f$ 也是有界闭区间。即对连续函数 $f:[a,b]\to\mathbb{R}$，存在 $\xi,\eta\in[a,b]$ 使得
\[
f(\xi)=\inf_{[a,b]}f,\qquad f(\eta)=\sup_{[a,b]}f,
\]
这就是极值定理（证明见 1021：$\operatorname{Im}f$ 有界且闭，故 $\sup\operatorname{Im}f$、$\inf\operatorname{Im}f$ 都属于 $\operatorname{Im}f$，即确界可以取到）。

\end{document}
